\section{\uppercase{Competitive Baselines}}
\label{sec:competitive}

\noindent
Experiment~3 compares the selected minimaps with methods that use the original movement signals and pretrained visual features. It uses a task set drawn from related work, which also shortens the handwriting protocol, and keeps all eight tasks as a secondary comparison.

\subsection{Task Set and Evaluation Protocol}
\noindent
The primary set comprises repeated \textit{l}, \textit{le} and \textit{les} (T2--T4), the tasks examined in the reference thesis~\cite{Casademunt2023}. We fix this set before fitting the models, and combining its three tasks by voting is our adaptation. Four methods are evaluated: static LPQ with SVM, LPQ with selection among the nine encodings, a kinematic/pressure SVM, and a frozen ResNet18 with a trained linear classification head.

The methods share the 25 outer participant partitions and three common inner folds. For each task and eligible encoding, mean inner-fold macro F1 selects the classifier parameters, and the two image methods that allow encoding selection choose one encoding for the task set by averaging these scores across its tasks, with ties resolved in favour of the static trace. Each method makes its own choices using only outer-training participants; standardization is fitted within the training partition, and selected models are refitted on the complete outer-training set. The three SVM methods retain the same 21-candidate linear/RBF search.

The primary endpoint is participant-level macro F1 from majority voting over T2--T4 (Section~\ref{sec:fusion}), with no calibration, learned weights or rule selection. Recording-level macro F1 over the three tasks is secondary. Holm adjustment covers the six method pairs separately at each endpoint, and the four within-method contrasts between three-task and eight-task voting form another family.

\subsection{Kinematics and Pressure}
\noindent
The signal baseline retains quantities removed by image fitting and per-record colour normalization. Its 83 fixed attributes, motivated by earlier kinematic and pressure measurements on PaHaW~\cite{Drotar2016}, comprise 11 acquisition times, contact fractions, stroke counts, path lengths and contact dimensions; 24 contact and in-air stroke duration and length summaries; 36 contact and in-air speed, acceleration and jerk summaries; and 12 contact pressure and absolute pressure-rate summaries. Each distribution contributes its mean, population standard deviation, median, 5th and 95th percentiles, and interquartile range. This is an independently specified descriptor, with no supervised feature screening.

Coordinates and pressure remain in native tablet units; timestamp differences are divided by 1,000 to obtain seconds. Velocity vectors are secants at interval midpoints. Successive vector differences over midpoint time differences yield acceleration and jerk, whose magnitudes are summarized. Derivatives never cross pen lifts or invalid intervals. No smoothing is applied, so quantization can influence higher derivatives.

A pre-fit audit excluded 255 of 1,364,022 intervals: 253 negative increments and two positive jumps exceeding 60s. None occurred within contact. Positive intervals up to 60s are retained, including 145 gaps of 1--26.61s; their velocities describe average movement across the gap. Acquisition time sums valid intervals, dividing each state-transition interval equally between contact and air. Stroke lengths and derivatives exclude transitions. Six zeros represent an absent distribution, recorded separately in diagnostics: 56 recordings lack in-air distributions and 57 lack an in-air jerk distribution. These conventions preserve every recording while making incomplete temporal information explicit.

\subsection{Transfer from a Pretrained CNN}
\noindent
ResNet18~\cite{He2016} supplies 512 features from its global average-pooling layer using the fixed torchvision \texttt{IMAGENET1K\_V1} weights. Images undergo the same foreground crop and square white padding as LPQ, followed by bilinear resizing to $224\times224$ and fixed ImageNet channel normalization. The backbone remains in evaluation mode, with convolutional weights and batch-normalization statistics frozen. There is no augmentation or end-to-end fine-tuning.

A class-balanced, $L_2$-regularized logistic head is fitted to standardized features. Its inner search evaluates seven values, $C\in\{10^{-4},10^{-3},10^{-2},10^{-1},1,10,100\}$, and selects among the same nine image encodings as LPQ. The comparison therefore concerns complete learning procedures with different feature families and search budgets.

\subsection{Comparison on the Three Fixed Tasks}
\begin{table*}[t]
\centering\small
\caption{Experiment~3 on T2--T4, fixed before fitting. Recording-level and participant-level macro F1 are distinct endpoints; the interval is the 95\% participant-bootstrap interval for participant-level macro F1, and each value averages five repetitions. Enc.~\# is the encoding of Table~\ref{tab:encoding} chosen most often across the 25 outer training sets; the full counts are 1$\times$2, 3$\times$4, 4$\times$1, 6$\times$4, 7$\times$4, 8$\times$2 and 9$\times$8 for selected LPQ, and 1$\times$1, 2$\times$4, 3$\times$1, 4$\times$3, 5$\times$2 and 7$\times$14 for ResNet18. Bold marks the highest value in each column.}
\label{tab:competitive}
\setlength{\tabcolsep}{4pt}
\begin{tabular}{@{}lcrrrr@{}}
\toprule
Method & Enc.~\# & Recording-level F1 & Participant-level F1 & 95\% CI & Accuracy (\%)\\
\midrule
Static LPQ + SVM & 1 & 0.5815 & 0.6143 & [0.5194, 0.7049] & 61.60\\
Training-selected LPQ + SVM & 9 (8/25) & 0.5773 & \textbf{0.6449} & [0.5502, 0.7340] & \textbf{64.80}\\
Kinematic/pressure SVM & -- & \textbf{0.6074} & 0.6279 & [0.5409, 0.7106] & 62.93\\
ResNet18 transfer + linear head & 7 (14/25) & 0.5710 & 0.6021 & [0.5136, 0.6873] & 60.27\\
\bottomrule
\end{tabular}
\end{table*}
\noindent
The training-selected minimap pipeline gives the best participant-level decisions of the four methods (Table~\ref{tab:competitive}). It has the highest macro F1 (0.6449), accuracy (64.80\%), balanced accuracy (0.6472), PD sensitivity (0.5892) and H specificity (0.7053). Its margin over static images is 0.0307 [$-$0.0107, 0.0748]: encoding the dynamics adds three points of macro F1 with the same descriptor and classifier, although a cohort of 75 participants cannot resolve margins of this size. Selection chose a dynamic encoding in 23 of the 25 training sets, most often $PAZ$ (9), whereas ResNet18 favoured $SPZ$ (7).

The kinematic descriptor leads the recording-level endpoint (0.6074) and the vote-fraction AUC (0.6691 against 0.6139), so the two pipelines are strongest on different endpoints. An average of per-task scores does not capture how errors coincide across tasks, which is what the participant-level vote rewards.

\subsection{Effect of the Task Subset}
\noindent
With all eight tasks, majority voting reaches participant-level macro F1 of 0.6330 for static LPQ, 0.6166 for selected LPQ, 0.5941 for kinematics and 0.6225 for ResNet18 features. Experiment~3 always votes, whereas Experiment~2 chose the rule within training, so the participant-level scores of the two experiments are not directly comparable.

The shorter protocol suits the minimap pipeline. Moving from eight tasks to T2--T4 raises its participant-level macro F1 by 0.0283 [$-$0.0358, 0.0950] and that of the kinematic descriptor by 0.0338, while static LPQ changes by $-$0.0188 and ResNet18 features by $-$0.0204. Recording-level means over three and eight tasks describe different task populations and are not compared directly.

\begin{table*}[t]
\centering\footnotesize
\caption{Recording-level macro F1 per task for the four methods of Experiment~3 evaluated on all eight tasks, averaged over five repetitions. Bold marks the highest value for each task and for the mean.}
\label{tab:per-task}
\setlength{\tabcolsep}{5pt}
\begin{tabular}{@{}lrrrrrrrrr@{}}
\toprule
Method & T1 & T2 & T3 & T4 & T5 & T6 & T7 & T8 & Mean\\
\midrule
Static LPQ + SVM & 0.5846 & 0.5505 & 0.5860 & 0.6079 & 0.5588 & 0.5232 & 0.5226 & 0.5620 & 0.5620\\
Training-selected LPQ + SVM & \textbf{0.6327} & 0.5407 & 0.5614 & 0.6042 & 0.5340 & 0.5362 & 0.5077 & 0.5772 & 0.5618\\
Kinematic/pressure SVM & 0.5074 & \textbf{0.6381} & \textbf{0.6662} & 0.5177 & 0.5062 & 0.5491 & 0.4773 & \textbf{0.6788} & 0.5676\\
ResNet18 transfer + linear head & 0.5984 & 0.5345 & 0.5457 & \textbf{0.6246} & \textbf{0.5772} & \textbf{0.5917} & \textbf{0.5888} & 0.5892 & \textbf{0.5813}\\
\bottomrule
\end{tabular}
\end{table*}

The eight tasks also show that the methods are complementary (Table~\ref{tab:per-task}). Selected minimaps give the highest recording-level macro F1 on the spiral, 0.6327 against 0.5846 for static LPQ, 0.5984 for ResNet18 and 0.5074 for kinematics, whereas kinematics lead T2, T3 and T8 and ResNet18 features lead T4--T7. Across all eight tasks, ResNet18 has the highest recording-level mean (0.5813, against 0.5618 for selected LPQ), which makes pretrained CNN features a competitive alternative classifier for the same minimaps.
